% Short summary deck: F-1 in NVIDIA OpenSMA, found by ESBMC, fixed upstream.
% Companion to fv_f1_executive.tex; intended for a quick scan (LinkedIn / X /
% blog link). 4 slides, 16:9.
%
% Build with:  pdflatex fv_f1_summary.tex  (run twice)
%
\documentclass[aspectratio=169,11pt]{beamer}

\usepackage[T1]{fontenc}
\usepackage{lmodern}
\usepackage{microtype}
\usepackage{xcolor}
\usepackage{booktabs}
\usepackage{listings}
\usepackage{pifont}

%% ---------- Theme ----------
\usetheme{Boadilla}
\useinnertheme{circles}
\useoutertheme{infolines}

% NVIDIA-inspired palette
\definecolor{nvgreen}{HTML}{76B900}
\definecolor{nvgreendark}{HTML}{4F8000}
\definecolor{nvblack}{HTML}{1A1A1A}
\definecolor{nvgray}{HTML}{4D4D4D}
\definecolor{nvlight}{HTML}{F2F2F2}
\definecolor{nvred}{HTML}{C0392B}
\definecolor{nvblue}{HTML}{2C3E50}

\setbeamercolor{structure}{fg=nvgreendark}
\setbeamercolor{title}{fg=nvblack}
\setbeamercolor{frametitle}{fg=white,bg=nvgreendark}
\setbeamercolor{author}{fg=nvgray}
\setbeamercolor{institute}{fg=nvgray}
\setbeamercolor{date}{fg=nvgray}
\setbeamercolor{block title}{fg=white,bg=nvgreendark}
\setbeamercolor{block body}{bg=nvlight}
\setbeamercolor{itemize item}{fg=nvgreendark}
\setbeamercolor{itemize subitem}{fg=nvgreendark}
\setbeamercolor{section in head/foot}{bg=nvblack,fg=white}
\setbeamercolor{author in head/foot}{bg=nvgreendark,fg=white}
\setbeamercolor{title in head/foot}{bg=nvblack,fg=white}
\setbeamercolor{date in head/foot}{bg=nvgreendark,fg=white}

\setbeamertemplate{navigation symbols}{}
\hypersetup{colorlinks=true,urlcolor=nvgreendark,linkcolor=nvgreendark}
\setbeamertemplate{itemize item}{\textbullet}
\setbeamerfont{frametitle}{series=\bfseries}
\setbeamerfont{title}{series=\bfseries,size=\LARGE}

%% ---------- Code listing style ----------
\lstdefinestyle{cppstyle}{
  language=C++,
  basicstyle=\ttfamily\scriptsize,
  keywordstyle=\color{nvgreendark}\bfseries,
  commentstyle=\color{nvgray}\itshape,
  stringstyle=\color{nvblue},
  numbers=none,
  showstringspaces=false,
  breaklines=true,
  frame=single,
  rulecolor=\color{nvgray!30},
  backgroundcolor=\color{nvlight},
  xleftmargin=0.4em,
  xrightmargin=0.4em,
  aboveskip=0.4em,
  belowskip=0.4em,
}
\lstset{style=cppstyle}

%% ---------- Helper macros ----------
\newcommand{\cmark}{\textcolor{nvgreendark}{\ding{51}}}
\newcommand{\xmark}{\textcolor{nvred}{\ding{55}}}
\newcommand{\hil}[1]{\textbf{\textcolor{nvgreendark}{#1}}}
\newcommand{\bad}[1]{\textbf{\textcolor{nvred}{#1}}}

%% ---------- Title metadata ----------
\title[FV in NVIDIA OpenSMA]{Formal Verification finds and fixes\\an MCTP control-path bug in NVIDIA OpenSMA}
\subtitle{ESBMC counterexample, confirmed and fixed upstream}
\author[Cordeiro]{Lucas C. Cordeiro}
\institute[Univ.\ of Manchester]{University of Manchester \\ \textit{lucas.cordeiro@manchester.ac.uk}}
\date{May 2026}

\begin{document}

%% =================================================================
\begin{frame}
  \titlepage
  \vspace{-1.0em}
  \begin{center}
    \scriptsize\textcolor{nvgray}{Full report: \texttt{verification/REPORT.md}
    \textbar{} Issue:
    \href{https://github.com/NVIDIA/OpenSMA/issues/1}{NVIDIA/OpenSMA\#1}}
  \end{center}
\end{frame}

%% =================================================================
\begin{frame}[fragile]{The bug --- a four-line view}
  \begin{columns}[T,onlytextwidth]
    \begin{column}{0.52\textwidth}
\begin{lstlisting}[basicstyle=\ttfamily\tiny]
// pdk-mctp-platforms-router-plat.cpp:34
void set_cur_eid(RoutingTable& rt,
                 Packet::InterfaceType iface,
                 uint8_t eid)
{
    rt.ec.cur_eid.at(iface) = eid;  // (*)
}
\end{lstlisting}

{\tiny\textcolor{nvgray}{Source: OpenSMA, \textcopyright{} NVIDIA Corporation,
Apache-2.0.}}

\vspace{0.4em}
\scriptsize
\textbf{The asymmetric guard}
\begin{itemize}\setlength\itemsep{0.1em}
  \item \texttt{cur\_eid} is sized \texttt{UsEnd = 2}.
  \item Validator rejects only \texttt{iface $\geq$ End = 18}.
  \item Any \texttt{iface $\in [2, 17]$} \hil{passes validation},
        then aborts at (\textasteriskcentered).
  \item Sibling \texttt{get\_cur\_eid()} \emph{has} the guard;
        \bad{\texttt{set\_cur\_eid()} did not.}
\end{itemize}
    \end{column}

    \begin{column}{0.46\textwidth}
      \begin{block}{Behaviour under production flags}
        \scriptsize
        With \texttt{-fno-exceptions -fno-rtti},
        \texttt{std::array::at(OOB)} calls \hil{\texttt{abort()}} ---
        the firmware terminates (SIGABRT, exit 134).
      \end{block}

      \vspace{0.3em}
      \begin{block}{Trigger}
        \scriptsize
        An MCTP Control SetEpId Request whose
        \texttt{packet\_interface} falls in the gap drives the
        production dispatch path into the OOB access.
      \end{block}
    \end{column}
  \end{columns}
\end{frame}

%% =================================================================
\begin{frame}[fragile]{What ESBMC proved --- in three states}
  \begin{columns}[T,onlytextwidth]
    \begin{column}{0.46\textwidth}
      \footnotesize
      \textbf{Setup}
      \begin{itemize}\setlength\itemsep{0.1em}
        \scriptsize
        \item Harness builds a Control SetEpId Request; 9 header
              fields constrained to pass \texttt{Validator::validate()}.
        \item \texttt{iface\_val} left \emph{nondeterministic} in $[2, 18)$.
        \item Calls production \texttt{on\_set\_endpoint\_id()} unmodified.
      \end{itemize}

      \vspace{0.4em}
      \textbf{Solver result}
      \begin{itemize}\setlength\itemsep{0.1em}
        \scriptsize
        \item 275 verification conditions; Bitwuzla returns SAT in
              \hil{$<$ 0.01\,s}.
      \end{itemize}

      \vspace{0.4em}
      \textbf{Production confirmation}
      \begin{itemize}\setlength\itemsep{0.1em}
        \scriptsize
        \item ESBMC-generated test case compiled and run natively
              $\Rightarrow$ \hil{SIGABRT (signal 6)}.
      \end{itemize}
    \end{column}

    \begin{column}{0.52\textwidth}
\begin{lstlisting}[basicstyle=\ttfamily\tiny,language={}]
ESBMC counterexample (abridged):

  State 5   iface_val = 2
            (passes assume >= UsEnd && < End)

  State 9   valid = 1
            (Validator::validate() returns true)

  State 10  Violated property:
            std::array::at out of range
            [i::0 < 2 fails]
            at set_cur_eid()
            -> cur_eid.at(2) on size-2 array

VERIFICATION FAILED
\end{lstlisting}

      \vspace{0.2em}
      {\scriptsize End-to-end proof from validator entry through
      dispatch to OOB write --- \hil{no code modification, no manual
      test input.}}
    \end{column}
  \end{columns}
\end{frame}

%% =================================================================
\begin{frame}{Outcome --- confirmed and fixed upstream}
  \begin{block}{Disclosure timeline}
    \footnotesize
    \begin{itemize}\setlength\itemsep{0.15em}
      \item Reported to NVIDIA as
            \href{https://github.com/NVIDIA/OpenSMA/issues/1}{\hil{NVIDIA/OpenSMA\#1}}
            with the ESBMC counterexample and a one-line fix.
      \item NVIDIA's OpenSMA team acknowledged the report and tracked the
            fix internally.
      \item \hil{Fix applied upstream} (2026-05-11): the
            \texttt{interface\,$\geq$\,UsEnd} bound check is now mirrored
            into \texttt{set\_cur\_eid()}, closing the asymmetric-guard gap.
    \end{itemize}
  \end{block}

  \vspace{0.3em}
  \begin{block}{Scope and caveats}
    \scriptsize
    The verification covered 22 selected modules of OpenSMA
    ($\sim$7{,}900\,LOC of first-party C++); not the full codebase, and
    not vendored components (\texttt{libspdm}, FreeRTOS, MCU SDK glue).
    Findings beyond F-1 are documented in the full report; some are
    latent (no current packet-driven path) and recorded as such.
    \hil{No claims of broad firmware insecurity are made or implied.}
  \end{block}

  \vspace{0.3em}
  \centering
  \footnotesize
  Full report and reproducers:
  \href{https://github.com/lucasccordeiro/NVIDIA-OpenSMA}{\texttt{github.com/lucasccordeiro/NVIDIA-OpenSMA}}
  \\[0.2em]
  ESBMC: \href{https://github.com/esbmc/esbmc}{\texttt{github.com/esbmc/esbmc}}

  \vspace{0.5em}
  {\tiny\textcolor{nvgray}{Independent research, University of Manchester.
  NVIDIA, OpenSMA, and related marks are trademarks of NVIDIA Corporation;
  no affiliation or endorsement implied.}}
\end{frame}

\end{document}
